% GENERATED FILE. Edit canonical lessons, then run npm run generate:books.
% canonical-source-hash: fnv1a64:fee5b87b28908096
% canonical-lessons: HI-C301-sayran, HI-C301-harbarana, HI-C301-bujhna, HI-C301-pighalna, HI-C301-phaphola

\chapter{Putting the Fire Out}
\label{ch:hi-putting-the-fire-out}
\hlchaptermodality{301}

\begin{chapteropening}
\textbf{By the end of this chapter:} \emph{I can say how a fire was put out.}

Complete the last lesson of chapter 301: I can say how a fire was put out.
\end{chapteropening}

\section[\texorpdfstring{\hi{सायरन} (sāyran)}{सायरन (sāyran)}]{\hi{सायरन} (sāyran) — a siren}

\label{lesson:HI-C301-sayran}



\begin{quote}
Before the new one: say the Hindi for light, then the Hindi for scorching heat.
\end{quote}

\subsection*{You'll want to know: \hi{सायरन}}
\textbf{\hi{सायरन}} — \emph{sāyran} — \textquotedblleft{}a siren\textquotedblright{}.

\begin{quote}
\hi{सड़क} \hi{पर} \hi{सायरन} \hi{बज} \hi{रहा} \hi{था।}
\end{quote}

\begin{quote}
\emph{saṛak par sāyran baj rahā thā.}
\end{quote}

\begin{quote}
A siren was sounding on the road.
\end{quote}

Masculine.

\begin{cousinweb}
\textbf{\hi{सायरन}} (\emph{sāyran}) is borrowed from English \emph{siren}, taken into Hindi as a masculine noun. Like a bell, a siren \textbf{\hi{बजता} \hi{है}}.
\end{cousinweb}

\begin{grammarlens}[title={What you've built}]
One more word for telling what happened.
\end{grammarlens}

\subsection*{Your turn}
\begin{itemize}
\raggedright
  \item \emph{Say it:} \emph{sāyran}
  \item \emph{Say it:} \emph{sāyran}, once more
  \item \emph{Say it:} \emph{saṛak par sāyran baj rahā thā.} — A siren was sounding on the road.
  \item \emph{Recall:} say the Hindi for light, then the Hindi for scorching heat, then say \emph{sāyran} again
\end{itemize}

\begin{tcolorbox}[breakable,colback=teal!4,colframe=teal!35!black,title={Before you move on}]
What does \hi{सायरन} mean? (\textquotedblleft{}A siren\textquotedblright{}.) Say it once more.
\end{tcolorbox}
\section[\texorpdfstring{\hi{हड़बड़ाना} (haṛbaṛānā)}{हड़बड़ाना (haṛbaṛānā)}]{\hi{हड़बड़ाना} (haṛbaṛānā) — to be flustered; to rush in a panic}

\label{lesson:HI-C301-harbarana}



\begin{quote}
Before the new one: say the Hindi for scorching heat, then the Hindi for a siren.
\end{quote}

\subsection*{You'll want to know: \hi{हड़बड़ाना}}
\textbf{\hi{हड़बड़ाना}} — \emph{haṛbaṛānā} — \textquotedblleft{}to be flustered; to rush in a panic\textquotedblright{}.

\begin{quote}
\hi{सायरन} \hi{सुनकर} \hi{लोग} \hi{हड़बड़ा} \hi{गए।}
\end{quote}

\begin{quote}
\emph{sāyran sunkar log haṛbaṛā gae.}
\end{quote}

\begin{quote}
Hearing the siren, people got flustered.
\end{quote}

Past (he …): \hi{वह} \hi{हड़बड़ा} \hi{गया}, \emph{vah haṛbaṛā gayā}.

\begin{cousinweb}
\textbf{\hi{हड़बड़ाना}} (\emph{haṛbaṛānā}) is set by CDIAL under \emph{*haḍabaḍa} (CDIAL 13949), a word CDIAL gives for confusion and very probably imitative. It is the rush when nobody knows what to do first.
\end{cousinweb}

\begin{grammarlens}[title={What you've built}]
One more word for telling what happened.
\end{grammarlens}

\subsection*{Your turn}
\begin{itemize}
\raggedright
  \item \emph{Say it:} \emph{haṛbaṛānā}
  \item \emph{Say it:} \emph{haṛbaṛānā}, once more
  \item \emph{Say it:} \emph{sāyran sunkar log haṛbaṛā gae.} — Hearing the siren, people got flustered.
  \item \emph{Recall:} say the Hindi for scorching heat, then the Hindi for a siren, then say \emph{haṛbaṛānā} again
\end{itemize}

\begin{tcolorbox}[breakable,colback=teal!4,colframe=teal!35!black,title={Before you move on}]
What does \hi{हड़बड़ाना} mean? (\textquotedblleft{}To be flustered; to rush in a panic\textquotedblright{}.) Say it once more.
\end{tcolorbox}
\section[\texorpdfstring{\hi{बुझना} (bujhnā)}{बुझना (bujhnā)}]{\hi{बुझना} (bujhnā) — to go out (of a fire or light)}

\label{lesson:HI-C301-bujhna}



\begin{quote}
Before the new one: say the Hindi for a siren, then the Hindi for to be flustered.
\end{quote}

\subsection*{You'll want to know: \hi{बुझना}}
\textbf{\hi{बुझना}} — \emph{bujhnā} — \textquotedblleft{}to go out (of a fire or light)\textquotedblright{}.

\begin{quote}
\hi{दो} \hi{घंटे} \hi{बाद} \hi{आग} \hi{बुझ} \hi{गई।}
\end{quote}

\begin{quote}
\emph{do ghaṇṭe bād āg bujh gaī.}
\end{quote}

\begin{quote}
Two hours later the fire went out.
\end{quote}

Past (it …): \hi{वह} \hi{आग} \hi{बुझ} \hi{गई}, \emph{vah āg bujh gaī}.

\begin{cousinweb}
\textbf{\hi{बुझना}} (\emph{bujhnā}) is inherited from \emph{*vijjhāyati}, a form CDIAL reconstructs (CDIAL 11703). It is the intransitive partner of \textbf{\hi{बुझाना}}, to put out: you \textbf{\hi{बुझाते} \hi{हैं}} a fire, and the fire \textbf{\hi{बुझती} \hi{है}}.
\end{cousinweb}

\begin{grammarlens}[title={What you've built}]
One more word for telling what happened.
\end{grammarlens}

\subsection*{Your turn}
\begin{itemize}
\raggedright
  \item \emph{Say it:} \emph{bujhnā}
  \item \emph{Say it:} \emph{bujhnā}, once more
  \item \emph{Say it:} \emph{do ghaṇṭe bād āg bujh gaī.} — Two hours later the fire went out.
  \item \emph{Recall:} say the Hindi for a siren, then the Hindi for to be flustered, then say \emph{bujhnā} again
\end{itemize}

\begin{tcolorbox}[breakable,colback=teal!4,colframe=teal!35!black,title={Before you move on}]
What does \hi{बुझना} mean? (\textquotedblleft{}To go out (of a fire or light)\textquotedblright{}.) Say it once more.
\end{tcolorbox}
\section[\texorpdfstring{\hi{पिघलना} (pighalnā)}{पिघलना (pighalnā)}]{\hi{पिघलना} (pighalnā) — to melt}

\label{lesson:HI-C301-pighalna}



\begin{quote}
Before the new one: say the Hindi for to be flustered, then the Hindi for to go out (of a fire or light).
\end{quote}

\subsection*{You'll want to know: \hi{पिघलना}}
\textbf{\hi{पिघलना}} — \emph{pighalnā} — \textquotedblleft{}to melt\textquotedblright{}.

\begin{quote}
\hi{आग} \hi{की} \hi{गर्मी} \hi{से} \hi{खिड़की} \hi{का} \hi{प्लास्टिक} \hi{पिघल} \hi{गया।}
\end{quote}

\begin{quote}
\emph{āg kī garmī se khiṛkī kā plāsṭik pighal gayā.}
\end{quote}

\begin{quote}
The heat of the fire melted the plastic on the window.
\end{quote}

Past (it …): \hi{वह} \hi{बर्फ़} \hi{पिघल} \hi{गई}, \emph{vah barf pighal gaī}.

\begin{cousinweb}
\textbf{\hi{पिघलना}} (\emph{pighalnā}) is listed by CDIAL under \emph{*praghārayati}, a reconstructed form (CDIAL 8485). Ice, wax and plastic \textbf{\hi{पिघलते} \hi{हैं}}; so, in speech, does a hard heart.
\end{cousinweb}

\begin{grammarlens}[title={What you've built}]
One more word for telling what happened.
\end{grammarlens}

\subsection*{Your turn}
\begin{itemize}
\raggedright
  \item \emph{Say it:} \emph{pighalnā}
  \item \emph{Say it:} \emph{pighalnā}, once more
  \item \emph{Say it:} \emph{āg kī garmī se khiṛkī kā plāsṭik pighal gayā.} — The heat of the fire melted the plastic on the window.
  \item \emph{Recall:} say the Hindi for to be flustered, then the Hindi for to go out (of a fire or light), then say \emph{pighalnā} again
\end{itemize}

\begin{tcolorbox}[breakable,colback=teal!4,colframe=teal!35!black,title={Before you move on}]
What does \hi{पिघलना} mean? (\textquotedblleft{}To melt\textquotedblright{}.) Say it once more.
\end{tcolorbox}
\section[\texorpdfstring{\hi{फफोला} (phapholā)}{फफोला (phapholā)}]{\hi{फफोला} (phapholā) — a blister}

\label{lesson:HI-C301-phaphola}



\begin{quote}
Before the new one: say the Hindi for to go out (of a fire or light), then the Hindi for to melt.
\end{quote}

\subsection*{You'll want to know: \hi{फफोला}}
\textbf{\hi{फफोला}} — \emph{phapholā} — \textquotedblleft{}a blister\textquotedblright{}.

\begin{quote}
\hi{आग} \hi{से} \hi{उसके} \hi{हाथ} \hi{पर} \hi{फफोले} \hi{पड़} \hi{गए।}
\end{quote}

\begin{quote}
\emph{āg se uske hāth par phaphole paṛ gae.}
\end{quote}

\begin{quote}
The fire left blisters on his hand.
\end{quote}

Masculine.

\begin{cousinweb}
\textbf{\hi{फफोला}} (\emph{phapholā}) is inherited from \emph{*prasphōṭa}, a form CDIAL reconstructs (CDIAL 8883), which gathers words for bubbles and blisters. Blisters \textbf{\hi{पड़ते} \hi{हैं}}: they come up on the skin.
\end{cousinweb}

\begin{grammarlens}[title={What you've built}]
One more word for telling what happened.
\end{grammarlens}

\subsection*{Your turn}
\begin{itemize}
\raggedright
  \item \emph{Say it:} \emph{phapholā}
  \item \emph{Say it:} \emph{phapholā}, once more
  \item \emph{Say it:} \emph{āg se uske hāth par phaphole paṛ gae.} — The fire left blisters on his hand.
  \item \emph{Recall:} say the Hindi for to go out (of a fire or light), then the Hindi for to melt, then say \emph{phapholā} again
\end{itemize}

\begin{tcolorbox}[breakable,colback=teal!4,colframe=teal!35!black,title={Before you move on}]
What does \hi{फफोला} mean? (\textquotedblleft{}A blister\textquotedblright{}.) Say it once more.
\end{tcolorbox}
